% ============================================================================
%  Tier-3 report: Cumulative Ecological Controls + Prioritized Improvement Plan
%  Self-contained: compile directly with  pdflatex  (no external figures).
%  Style follows 22_6_Continuous_Observation_New_Baselines.tex.
%  Purpose:
%   (i)   recap the Tier-2 continuous-state, noisy-observation POMDP and what the
%         executed run showed (so this file stands alone);
%   (ii)  diagnose why the Tier-2 action/reward design is a bottleneck;
%   (iii) specify the redesigned setting: model-only, *cumulative*, ecologically
%         *capped* controls on (r,K), with revised 5- and 10-action tables;
%   (iv)  give the prioritized list of improvements worth doing to recover a
%         learned advantage over PLUS and MOOR (especially at high noise);
%   (v)   give an implementation checklist for the coding agent that built the
%         Tier-2 repo, mapping each change to the code it must touch.
% ============================================================================
\documentclass[11pt]{article}
\usepackage[utf8]{inputenc}
\usepackage[margin=0.9in]{geometry}
\usepackage{amsmath,amssymb}
\usepackage{booktabs}
\usepackage{longtable}
\usepackage{enumitem}
\usepackage{xcolor}
\usepackage[hidelinks]{hyperref}

\newcommand{\reff}{r_{\mathrm{eff}}}
\newcommand{\Keff}{K_{\mathrm{eff}}}
\newcommand{\rbase}{r_{\mathrm{base}}}
\newcommand{\Kbase}{K_{\mathrm{base}}}
\newcommand{\Kref}{K_{\mathrm{ref}}}
\newcommand{\rmax}{r_{\max}}
\newcommand{\rmin}{r_{\min}}
\newcommand{\Kmax}{K_{\max}}
\newcommand{\Kmin}{K_{\min}}
\newcommand{\sobs}{\sigma_{\mathrm{obs}}}
\newcommand{\sthr}{s_{\mathrm{safe}}}
\newcommand{\Rtrue}{R^{\mathrm{true}}}

\title{Adaptive Management under Observation Error (Tier-3):\\
Cumulative, Ecologically-Capped Controls and a Prioritized Improvement Plan\\
\large (from one-step parameter nudges to persistent management effort;
restoring a learned advantage over PLUS/MOOR)}
\author{Deep RL for Ecological Population Management}
\date{2026-06-27}

\begin{document}
\maketitle

\begin{abstract}
This report builds directly on the Tier-2 continuous-state, noisy-observation
benchmark (spec: \texttt{22\_6\_Continuous\_Observation\_New\_Baselines.tex};
executed results: \texttt{22\_6\_Experiment\_Results}) and is written to stand on
its own. We first recap, briefly, the Tier-2 setting and the headline empirical
finding: general dynamics-learning learners beat the stronger ecology baseline
$\max(\text{PLUS},\text{MOOR})$ in $33$--$83\%$ of cells at low-to-moderate survey
noise ($\sobs\!\le\!0.2$), but that advantage \emph{collapses to $\approx 0$} at
high noise ($\sobs{=}0.4$), where the single-model mechanistic baseline (MOOR) and
the guarded learner (OGSRL) become best. We then diagnose a design flaw that is
partly responsible: in the executed Tier-2 code, actions perturbed the growth
parameters $(r,K)$ for a \emph{single step} and then reset, the abundance-reward
saturated at a \emph{fixed} reference, and there was no direct authority over
abundance, with the combined effect that the ``conservation'' actions were
economically dominated and the decision-relevant signal was squeezed into a thin
near-threshold band where observation noise is most destructive. We redesign the
control interface accordingly: actions act \emph{only} through the population model
and their effect on $(r,K)$ is \emph{cumulative} (persistent management effort)
and \emph{capped} for ecological and dynamical realism (a species cannot be made
to grow faster than its biology, and the discrete maps must stay in their stable
regime). We give the augmented transition, the caps, the corrected reward, and
revised $5$- and $10$-action tables. Finally we give a prioritized list of
improvements (Tiers A--D) targeted at recovering a learned advantage under noise,
and a concrete implementation checklist for the coding agent that maps each change
onto the Tier-2 repository.
\end{abstract}

\section{Relationship to the Tier-2 report and what changed}
The Tier-2 spec posed the central question: \emph{under decision-relevant
misspecification and observation error, can general offline model-based learners
that infer the dynamics and filter the state beat the ecology-specific mechanistic
baselines (PLUS, MOOR) and avoid collapse?} The executed Tier-2 run answered
``yes at $\sobs\!\le\!0.2$, no at $\sobs{=}0.4$.'' This report keeps the question
and the baselines unchanged but makes two kinds of change:
\begin{itemize}[nosep,leftmargin=1.4em]
\item \textbf{A setting change (Section~\ref{sec:tier3}).} The control interface is
redesigned so management effort is \emph{cumulative} and acts only through the
model parameters $(r,K)$ with ecological caps. This fixes a concrete flaw
(Section~\ref{sec:diagnosis}) that made the Tier-2 conservation actions inert and
made the high-noise regime degenerate to a static-regulator problem.
\item \textbf{An improvement plan (Sections~\ref{sec:plan}--\ref{sec:checklist}).}
A prioritized set of method/diagnostic changes aimed squarely at the
$\sobs{=}0.4$ failure, plus a code-level checklist for the agent that implemented
Tier-2.
\end{itemize}
Everything needed to read this file is restated below; the Tier-2 documents are
referenced only for full numerical provenance.

\section{Recap: the continuous-state, noisy-observation setting}
\label{sec:recap}
The physical abundance is a continuous, unbounded, non-negative latent variable
$s_t\in\mathbb{R}_{\ge 0}$ with $s{=}0$ absorbing. A hidden per-episode parameter
vector $m=(\rbase,C,\theta,z_t,\dots)$ is drawn once per episode and never logged;
the complete latent state is $\zeta_t=(s_t,m)$, of which the agent observes neither
component directly. The manager observes only a noisy real-valued survey through a
multiplicative (log-normal) channel,
\begin{equation}
o_t = s_t\,\eta_t,\qquad \eta_t\sim\mathrm{LogNormal}(0,\sobs^2),
\label{eq:obs}
\end{equation}
so $\log o_t=\log s_t+\varepsilon_t$, $\varepsilon_t\sim\mathcal{N}(0,\sobs^2)$:
relative error is constant across abundance, $\sobs{=}0$ recovers a fully observed
MDP (still with hidden $m$), and increasing $\sobs$ smoothly increases partial
observability. The agent acts on the observation history
$h_t=(o_0,a_0,R_0,\dots,o_t)$ through $\pi(a_t\mid h_t)$ or $\pi(a_t\mid b_t)$ with
a belief $b_t\approx p(s_t,m\mid h_t)$.

The latent abundance evolves under a base Ricker map plus three single-structure
stress models, each adding one hidden per-episode structure:
\begin{align}
\textbf{Ricker (base):}\quad
& s_{t+1}=s_t\exp\!\big(\reff(1-s_t/\Keff)\big),\\
\textbf{Allee--Ricker } (C):\quad
& s_{t+1}=s_t\exp\!\big(\reff(1-s_t/\Keff)(s_t/C-1)\big),\\
\textbf{Theta-logistic } (\theta):\quad
& s_{t+1}=s_t+\reff\,s_t\big(1-(s_t/\Keff)^{\theta}\big),\\
\textbf{Regime-switch } (z_t):\quad
& s_{t+1}=s_t\exp\!\big(m_{z_t}\reff(1-s_t/\Keff)(s_t/C_{z_t}-1)\big),
\end{align}
where below $C$ the Allee factor $(s_t/C-1)$ flips sign (critical depensation).
The simulator uses these only to generate data; the learner must infer the law.

The action space is a finite set $\mathcal{A}=\{0,\dots,A-1\}$, $A\in\{5,10\}$.
\textbf{The action--dynamics interface is what this report changes}
(Sections~\ref{sec:diagnosis}--\ref{sec:tier3}). The reward, in Tier-2, was
computed from the observation to avoid leaking the true state into the dataset,
with a saturating abundance utility and a one-time collapse penalty,
\begin{equation}
R_t=\underbrace{\alpha\,\frac{o_t}{o_t+\Kref}}_{\text{abundance (from }o_t\text{)}}
-\mathrm{cost}(a_t)-P\cdot\mathbb{1}\big[s_t>\sthr\wedge s_{t+1}\le\sthr\big],
\qquad \alpha{=}1,\ \Kref{=}\Kbase{=}500,\ P{=}20,
\label{eq:reward}
\end{equation}
with $\sthr{=}50$ the safety floor and $P{=}20$ chosen so one collapse outweighs a
whole good episode (achievable return $\sim 7$--$9$). A true-state metric $\Rtrue$
is logged for evaluation only.

\section{What Tier-2 showed, and the one-line diagnosis}
\label{sec:results}
The executed matrix was
$\{\text{allee},\text{theta},\text{regime}\}\times\{5\mathrm{a},10\mathrm{a}\}
\times\sobs\!\in\!\{0,0.1,0.2,0.4\}\times 7$ methods, with a shared $1024$-particle
belief filter, $5$ held-out seed blocks ($250$ episodes/cell), horizon $50$,
$\gamma{=}0.95$. The two headline tables (means over the six family cells per
$\sobs$, learned filter) were:

\begin{longtable}{@{}lccccc@{}}
\caption{Operational return by method and noise level (Tier-2).}\\
\toprule
Method & $\sobs{=}0$ & $0.1$ & $0.2$ & $0.4$ & all-$\sobs$ \\
\midrule
\endfirsthead
\toprule
Method & $\sobs{=}0$ & $0.1$ & $0.2$ & $0.4$ & all-$\sobs$ \\
\midrule
\endhead
RefPlan & $7.33$ & $7.23$ & $7.14$ & $5.14$ & $6.71$ \\
MOPO    & $7.20$ & $7.21$ & $6.94$ & $5.77$ & $6.78$ \\
MOOR \emph{(baseline)} & $6.97$ & $6.94$ & $7.22$ & $\mathbf{7.80}$ & $\mathbf{7.23}$ \\
BA-MCTS & $6.01$ & $5.87$ & $5.66$ & $4.18$ & $5.43$ \\
PLUS \emph{(baseline)} & $5.64$ & $5.74$ & $5.93$ & $6.56$ & $5.97$ \\
OGSRL   & $4.37$ & $5.95$ & $7.39$ & $7.48$ & $6.30$ \\
Delphic & $5.58$ & $3.74$ & $5.98$ & $5.77$ & $5.27$ \\
\bottomrule
\end{longtable}

\begin{longtable}{@{}lcccc@{}}
\caption{Beats-both rate: fraction of $6$ cells beating $\max(\text{PLUS},\text{MOOR})$.}\\
\toprule
General learner & $\sobs{=}0$ & $0.1$ & $0.2$ & $0.4$ \\
\midrule
\endfirsthead
\toprule
General learner & $\sobs{=}0$ & $0.1$ & $0.2$ & $0.4$ \\
\midrule
\endhead
BA-MCTS & $0.67$ & $\mathbf{0.83}$ & $0.50$ & $0.00$ \\
RefPlan & $0.67$ & $0.50$ & $0.33$ & $0.00$ \\
MOPO    & $0.50$ & $0.33$ & $0.33$ & $0.00$ \\
OGSRL   & $0.00$ & $0.17$ & $0.50$ & $0.17$ \\
Delphic & $0.00$ & $0.17$ & $0.17$ & $0.00$ \\
\bottomrule
\end{longtable}

\noindent\textbf{The decisive axis is observation noise, and it reverses the
ranking.} At $\sobs{=}0$ the general learners lead; as $\sobs$ rises the learned
advantage erodes in lockstep with belief-filter RMSE (averaged over methods
$\approx 0,43,80,160$ at $\sobs{=}0,0.1,0.2,0.4$; at the high end comparable to the
state scale), and at $\sobs{=}0.4$ the beats-both rate is $\approx 0$. The
single-model MOOR is best overall ($7.23$) and \emph{improves} with noise. The
mechanism is: at high $\sobs$ the belief is too degraded to recover the per-episode
parameters, so a static conservative regulator (one Ricker fit, regulate to a safe
mid-band) is robust, while the belief-based planners propagate a degraded belief
through a model and a multi-step plan and compound the error. \textbf{Two things
follow:} (1) the learners' advantage is real but observability-limited; (2) part of
the high-noise outcome is an artifact of the \emph{control interface}, which is the
subject of Section~\ref{sec:diagnosis}.

\section{Why the Tier-2 action/reward design is a bottleneck}
\label{sec:diagnosis}
The executed Tier-2 action tables perturbed the growth parameters additively,
$\reff(a)=\rbase+\Delta r(a)$ and $\Keff(a)=\Kbase+\Delta K(a)$, with $\Delta r\in
\{-0.02,\dots,+0.02\}$ and $\Delta K\in\{0,+100,+200\}$, and (crucially) these
perturbations applied to a \emph{single} transition and then reset. This produced
four coupled problems:
\begin{enumerate}[leftmargin=1.6em,nosep]
\item \textbf{Restoration is economically dominated.} With a one-step $\Delta K$
and a reward that saturates at the fixed $\Kref{=}500$, paying e.g.\ $0.30$ for a
one-step $\Keff:500\!\to\!700$ buys $\le\!0.08$ of abundance reward that step, with
negligible carryover, and the bump evaporates next step. The actions that are
supposed to be the heart of adaptive management are the ones a rational agent
almost never takes.
\item \textbf{Harvest as $-\Delta r$ is biologically backwards (and, at one step,
weak).} Culling removes individuals; it does not lower the intrinsic per-capita
growth rate for one step and then restore it.
\item \textbf{No rescue lever.} With actions acting only on $(r,K)$ for one step,
and no direct authority over $s$, a stock that is already below the Allee threshold
$C$ cannot be pulled back: below $C$ the factor $(s/C-1)$ is negative regardless of
$r$ or $K$, and raising $r$ makes the decline \emph{faster}. So a low start below
$C$ is unrecoverable.
\item \textbf{The decision-relevant signal is squeezed where noise is worst.}
Because conservation is dominated and rescue is impossible, the only consequential
decisions live in a thin band just above $C$ (prevent crossing). That band is
exactly where the log-normal survey is least informative about position relative to
$C$, so at $\sobs{=}0.4$ the controllable signal is buried in noise --- which is a
coherent, design-level reason the learned advantage vanishes.
\end{enumerate}
The fix is not to make the learners fancier but to make the \emph{controls
consequential}: persistent (so investment compounds), acting through the model (so
inferring the hidden structure is necessary to act well), and capped (so the
problem stays ecologically and numerically sane).

\section{The redesigned setting (Tier-3): model-only, cumulative, capped controls}
\label{sec:tier3}

\paragraph{Design principle.} Actions act \emph{only} through the population
model's parameters $(r,K)$; there is no direct authority over $s$. Direct stocking
would be \emph{structure-agnostic} --- adding individuals works regardless of the
hidden dynamics --- and would let a manager avoid ever understanding the model,
weakening the very capability the study tests. Routing all effect through $(r,K)$
makes the action's consequence depend on the hidden structure (e.g.\ raising $r$
helps \emph{above} the Allee threshold but \emph{accelerates collapse below it}),
so a learner must infer the structure to act well, while a Ricker-only baseline
gets the sign wrong near the threshold. This sharpens, rather than dilutes, the
PLUS/MOOR comparison.

\subsection{Augmented, cumulative transition}
Introduce two \emph{known} state variables, the accumulated management effort on
growth rate $\rho_t$ and on capacity $\kappa_t$. Per step, in order:
\begin{align}
\rho_t &= \rho_{t-1}+\Delta r(a_t), \qquad \rho_0=0, \\
\kappa_t &= \kappa_{t-1}+\Delta K(a_t), \qquad \kappa_0=0, \\
r_{\mathrm{eff},t} &= \mathrm{clip}\big(\rbase+\rho_t,\ \rmin,\ \rmax\big), \\
K_{\mathrm{eff},t} &= \mathrm{clip}\big(\Kbase+\kappa_t,\ \Kmin,\ \Kmax\big), \\
s_{t+1} &= f_{\mathrm{model}}\big(s_t;\ r_{\mathrm{eff},t},\ K_{\mathrm{eff},t}\big),
\label{eq:cumdyn}
\end{align}
where $f_{\mathrm{model}}$ is the per-cell Ricker/Allee/theta/regime map. The
agent observes $o_{t+1}=s_{t+1}\eta_{t+1}$ as before. Because $\rho_t,\kappa_t$ are
the agent's own cumulative actions, $(r_{\mathrm{eff},t},K_{\mathrm{eff},t})$ are \emph{known
exactly}; only $s$ remains partially observed and only $m$ remains hidden, so the
inference problem is unchanged in spirit but the controls now persist.

\paragraph{Optional maintenance decay.} To avoid an ``invest-then-coast'' optimum
and to model capacity that erodes when unmanaged, one may replace the accumulators
with leaky versions $\rho_t=(1-d_r)\rho_{t-1}+\Delta r(a_t)$ (and likewise
$\kappa_t$). The default is $d_r=d_K=0$ (pure cumulative); revisit after
calibration.

\subsection{Ecological and dynamical caps}
The caps are what keep cumulative controls meaningful rather than letting a policy
drive $r$ or $K$ to absurd values.
\begin{longtable}{@{}p{2.6cm}p{4.0cm}p{7.4cm}@{}}
\toprule
Cap & Value & Rationale \\
\midrule
\endhead
$\rmax$ (Ricker, Allee, regime) & $0.55$ & $r$ is the intrinsic rate of increase,
set by life history; management can raise the \emph{realized} rate toward the
biological ceiling (remove mortality) but not past it. $0.55$ reads as best-case
intensive management above an $\rbase\!\le\!0.30$ prior, and is far below the
Ricker instability onset (period-doubling at $r{=}2$, chaos near $2.5$). \\
$\rmax(\theta)$ (theta cells) & $1.6/\theta$ & The theta-logistic map is stable at
$K$ only for the \emph{product} $r\theta<2$; linearizing $s_{t+1}=s+rs(1-(s/K)^
\theta)$ about $K$ gives multiplier $(1-r\theta)$. The cap must therefore scale
with $\theta$: $\approx 0.53$ at $\theta{=}3$, $\approx 0.27$ at $\theta{=}6$
($1.6$ leaves a margin below $2$). \\
$\rmin$ & $-0.15$ (tune) & Sustained over-exploitation can push \emph{net} growth
negative (harvest mortality exceeds recruitment), so aggressive harvest genuinely
declines and can collapse the stock instead of merely freezing it. This is what
gives harvest real teeth now that it acts through $r$ rather than removing
individuals. \\
$\Kmin,\Kmax$ & $500,\ 1500$ & $\Kmin=\Kbase$ (no degradation below baseline in the
default, no-decay design); $\Kmax$ is a finite landscape ceiling --- restoration
cannot expand habitat without bound. \\
\bottomrule
\end{longtable}

\subsection{Reward: keep a \emph{fixed} reference (do not track $\Keff$)}
\label{sec:rewardfix}
A natural instinct is to let the saturation reference track capacity so that
restoration ``pays.'' It does the opposite. With $R\propto o/(o+\Keff)$, at
ecological steady state the population fills its capacity ($s\approx\Keff$), so
$o/(o+\Keff)\approx 0.5$ \emph{regardless of} $\Keff$, and immediately after a
restoration step ($\Keff$ jumps, $s$ has not caught up) the reward \emph{drops}.
Tracking $\Keff$ thus normalizes the gain away and restoration becomes safety-only
again. The combination that makes restoration worth its cost is a \emph{fixed}
reference \emph{plus} the cumulative dynamics of Section~\ref{sec:tier3}: a
persistently higher $\Keff$ sustains a higher absolute abundance, which a fixed
reference rewards (with diminishing returns). We therefore keep
\begin{equation}
R_t=\alpha\,\frac{o_t}{o_t+\Kref}-\mathrm{cost}(a_t)
-P\cdot\mathbb{1}\big[s_t>\sthr\wedge s_{t+1}\le\sthr\big],
\qquad \alpha{=}1,\ \Kref{=}500\ \text{(fixed)},\ P{=}20.
\label{eq:reward3}
\end{equation}
(Quantitatively: building $\Keff$ toward $\sim\!1000$ sustains $o\sim\!1000$ and
raises the abundance term from $0.50$ to $\approx 0.67$, a persistent
$+0.17$/step that can clear an early restoration investment; pushing further to
$1500$ adds only $\approx +0.08$/step, so the optimum is interior, not
``max-out-$K$.'') See Section~\ref{sec:plan}, Tier~A, for the orthogonal and
recommended upgrade to a belief-expected reward.

\subsection{Harvest as cumulative exploitation pressure}
Each harvest action applies a negative $\Delta r$ that accumulates in $\rho_t$ and
is floored at $\rmin$. Sustained aggressive harvest therefore drives $\reff$ down,
slows or reverses growth, and --- near the Allee threshold or under regime stress
--- precipitates collapse; the per-step revenue (a negative cost) must be weighed
against the delayed loss of abundance and the one-time $-P$. This recreates the
classic over-exploitation tragedy as an emergent property of the cumulative
dynamics, which is the richer decision problem chosen for this study.

\subsection{Revised action tables}
The only structural numeric change versus Tier-2 is that $\Delta K$ shrinks
$\sim\!5\times$ (it now accumulates, so $+200$/step would blow past any sane $K$
within a few steps), and harvest revenue is trimmed slightly; the $\Delta r$
increments are unchanged because $\pm0.01/\pm0.02$ per step ramps $r$ toward its
cap over $\sim\!15$--$25$ steps, so the \emph{timing} of investment matters. All
magnitudes are starting values to be calibrated against the $[0.15,0.24]$ collapse
band and the decision-relevance gate (Section~\ref{sec:checklist}).

\begin{longtable}{@{}c p{3.0cm} r r r p{4.6cm}@{}}
\caption{Revised 5-action set (cumulative $\Delta r,\Delta K$; fixed $\Kref$).}\\
\toprule
$a$ & Action & $\Delta r$/step & $\Delta K$/step & Cost/step & Interpretation \\
\midrule
\endfirsthead
\toprule
$a$ & Action & $\Delta r$/step & $\Delta K$/step & Cost/step & Interpretation \\
\midrule
\endhead
0 & Do Nothing & $0$ & $0$ & $0.00$ & passive monitoring \\
1 & Aggressive Harvest & $-0.02$ & $0$ & $-0.30$ & exploitation mortality (revenue); $\reff$ floored at $\rmin$ \\
2 & Predator / Disease Control & $+0.01$ & $0$ & $+0.20$ & reduces mortality, raises realized $r$ \\
3 & Intensive Restoration & $0$ & $+40$ & $+0.30$ & reserve expansion / revegetation; raises $K$ \\
4 & Flagship Conservation & $+0.02$ & $+40$ & $+0.60$ & breeding $+$ intensive restoration (both channels) \\
\bottomrule
\end{longtable}

\begin{longtable}{@{}c p{3.2cm} r r r p{4.4cm}@{}}
\caption{Revised 10-action set (cumulative $\Delta r,\Delta K$; fixed $\Kref$).}\\
\toprule
$a$ & Action & $\Delta r$/step & $\Delta K$/step & Cost/step & Interpretation \\
\midrule
\endfirsthead
\toprule
$a$ & Action & $\Delta r$/step & $\Delta K$/step & Cost/step & Interpretation \\
\midrule
\endhead
0 & Do Nothing & $0$ & $0$ & $0.00$ & passive monitoring \\
1 & Sustainable Harvest & $-0.01$ & $0$ & $-0.15$ & light exploitation pressure (revenue) \\
2 & Aggressive Harvest & $-0.02$ & $0$ & $-0.30$ & heavy exploitation pressure (revenue) \\
3 & Predator / Disease Control & $+0.01$ & $0$ & $+0.20$ & mortality reduction \\
4 & Breeding / Recruitment Support & $+0.02$ & $0$ & $+0.40$ & captive breeding, supplementary feeding \\
5 & Moderate Restoration & $0$ & $+20$ & $+0.15$ & fencing, fire management, planting \\
6 & Intensive Restoration & $0$ & $+40$ & $+0.30$ & reserve expansion, revegetation \\
7 & Integrated Conservation (light) & $+0.01$ & $+20$ & $+0.35$ & predator control $+$ moderate restoration \\
8 & Adaptive Conservation Trial & $+0.01$ & $+40$ & $+0.50$ & predator control $+$ intensive restoration \\
9 & Flagship Conservation Programme & $+0.02$ & $+40$ & $+0.60$ & breeding $+$ intensive restoration \\
\bottomrule
\end{longtable}

\subsection{Two companion fixes this design forces}
\begin{enumerate}[leftmargin=1.6em,nosep]
\item \textbf{The theta base prior is already unstable in its high corner.} With
$\rbase\sim\mathcal{U}(0.18,0.40)$ and $\theta\sim\mathcal{U}(3,6)$, the product
$r\theta$ reaches $2.4$ --- past the $r\theta{=}2$ stability boundary --- before
any action; this is almost certainly the ``high-$r/\theta$ overshoots and crashes''
behavior noted for theta in Tier-2. Draw with headroom under the cap, e.g.\ draw
$\theta$ first, then $\rbase\sim\mathcal{U}\big(0.12,\min(0.40,\,1.2/\theta)\big)$,
so base $r\theta\le 1.2$ and the action cap holds $r\theta<1.6$.
\item \textbf{Collapse accounting is prevention-only.} With no direct stocking, a
start (or fall) below $C$ is unrecoverable. So either initialise low-start episodes
just \emph{above} $C$ (rescuable-by-prevention), or report the collapse rate
\emph{split by start condition}, so the metric reflects policy skill rather than
start luck.
\end{enumerate}

\section{Prioritized improvement plan --- what's worth doing}
\label{sec:plan}
Ordered by expected impact on the central question (recovering a learned advantage
over $\max(\text{PLUS},\text{MOOR})$, especially at $\sobs{=}0.4$) and on
faithfulness. Tags refer to the user's six points and the Tier-2 report's eight AI
proposals.

\paragraph{Tier A --- directly attacks why learners lose at $\sobs{=}0.4$ (do these).}
\begin{enumerate}[leftmargin=1.8em,nosep]
\item \textbf{Joint state $+$ parameter belief filter} (Rao--Blackwellized particle
filter, or a recurrent / state-space encoder, over $(s,m)$). The Tier-2 filter
estimated only $s$; the per-episode parameters were a \emph{generic, unidentified}
latent, so the Bayes-adaptive edge that is the whole thesis was effectively
disabled. Particle \emph{count} ($1024$) is not the bottleneck for a low-dimensional
state; identifying $(r,C,\theta,z)$ is. \emph{Highest expected impact.} [user~\#1,
AI~\#1/\#2]
\item \textbf{Belief-expected reward} $r(b_t,a)=\mathbb{E}_{s\sim b_t}\!\big[\alpha
s/(s+\Kref)-\mathrm{cost}(a)-P\cdot\Pr(\text{collapse})\big]$ instead of the
single-observation reward; evaluate on $\Rtrue$. This is the POMDP-correct
belief-MDP reward, does not leak $s$ (it uses only the observation-derived belief),
and removes the single-sample noise that the multi-step planners \emph{compound}
through their rollouts --- a variance source MOOR's static fit ignores, so the fix
differentially helps the learners. [user~\#2, AI~\#7]
\item \textbf{Well-posed reward/action economics} --- realized by the Tier-3
control redesign (Section~\ref{sec:tier3}): cumulative, capped $(r,K)$ with a fixed
reference, harvest as exploitation pressure, and a real interior optimum for
restoration. Makes adaptive timing pay and removes the dominated-action regime.
[user~\#3]
\end{enumerate}

\paragraph{Tier B --- diagnostics that reinterpret the headline (cheap, do alongside A).}
\begin{enumerate}[leftmargin=1.8em,nosep]
\setcounter{enumi}{3}
\item \textbf{Oracle-filter ablation $+$ decision-relevance gate at $\sobs{=}0.4$.}
Feed the true $s$ to all methods at $\sobs{=}0.4$: if the learners then beat MOOR,
the bottleneck is the belief, not the algorithm. Run the oracle-MPC-vs-Ricker-MPC
gate \emph{at} $\sobs{=}0.4$ (Tier-2 enforced it only at $\sobs\!\le\!0.2$): if the
gap falls below the $1.0$ reward / $0.05$ collapse thresholds, the regime is no
longer decision-relevant and ``MOOR wins'' is the expected, honest outcome rather
than a learner failure. [user~\#6, AI implicitly]
\item \textbf{Report the pure-mechanistic MOOR ($6.83$, its own Ricker filter) and
$\Rtrue$ throughout.} The headline $7.23$ for MOOR benefits $+0.40$ from the shared
learned front-end; the honest baseline is lower. $\Rtrue$ is required for any
cross-$\sobs$ claim. [user~\#6, AI~\#7]
\end{enumerate}

\paragraph{Tier C --- per-method fixes (each can flip a weak method to competitive).}
\begin{enumerate}[leftmargin=1.8em,nosep]
\setcounter{enumi}{5}
\item \textbf{Delphic $\rightarrow$ model-based}, with a world-model-variance
penalty on the \emph{shared} MPC backbone: an ensemble of compatible world models
over the hidden-confounder ambiguity (ideally samples from the $m$-posterior of the
Tier-A filter), planned with the same MPC as MOPO/RefPlan/BA-MCTS, penalized by the
cross-world variance of value/return ($\approx$ a $Q$-penalty). Replaces the
underperforming linear-CQL head and makes the comparison clean (same backbone,
different uncertainty source). Validate that $u_\Delta$ increases monotonically
with $\sobs$. [user~\#4, AI~\#3]
\item \textbf{OGSRL cumulative-budget recalibration $+$ per-noise tuning.} The
Tier-2 deployment limit reused the discounted-cumulative budget ($0.02$) as a
\emph{one-step} probability threshold --- a scale mismatch ($\sum_t\gamma^t\approx
18.5$ over the horizon, so $0.02$ cumulative implies a per-step probability of
$\sim\!0.001$). State the semantics explicitly, set the cumulative budget from the
calibration target, decouple the deployment one-step threshold, and recalibrate per
$\sobs$. Likely recovers the poor $\sobs{=}0$ return ($4.37$) and could make OGSRL
the best \emph{safe} learner across all $\sobs$. [user~\#5, AI~\#4]
\item \textbf{Richer dynamics member $+$ BA-MCTS safety-aware leaves.} Add a small
MLP ensemble member (still composed with the exact emission) to help all
dynamics-learning methods where the belief is noisy; replace BA-MCTS's zero
terminal values with a fitted belief-value plus a near-threshold penalty to curb
its high collapse rate. Incremental. [AI~\#2, \#5]
\end{enumerate}

\paragraph{Tier D --- credibility.}
\begin{enumerate}[leftmargin=1.8em,nosep]
\setcounter{enumi}{8}
\item \textbf{More seeds / independent datasets per cell $+$ paired-bootstrap CIs.}
Beats-both over $6$ cells/level gives very wide intervals; add independently
trained datasets per cell (a robustness study) and report paired bootstrap
intervals as the headline, not point rates. [AI~\#6]
\end{enumerate}

\paragraph{If only three.} Tier~A items \#1 (filter with parameter identification),
\#2 (belief-expected reward), and Tier~B \#4 (oracle-filter $+$ gate at
$\sobs{=}0.4$): the first two are the most likely levers to push beats-both above
$0$ at high noise, and the third tells you whether that is even attainable or
whether $\sobs{=}0.4$ is genuinely un-adaptive.

\section{Implementation checklist for the coding agent}
\label{sec:checklist}
This section is written for the agent that implemented the Tier-2 repository
(\texttt{discrete\_action\_cont\_obser/}: environment, trajectory-preserving
collector with calibration profiles, shared $1024$-particle filter with a learned
log-space proposal, a $15$-member per-action $[1,x,x^2]$ ridge dynamics ensemble
with a receding-horizon MPC, the seven methods, and a unified evaluator). Each item
notes the Tier it serves. Preserve the unified evaluator, seed protocol, and the
public/private schema guard throughout.

\paragraph{(C0) Environment --- the Tier-3 control interface (setting change).}
\begin{enumerate}[leftmargin=1.8em,nosep]
\item Add two persistent, \emph{known} state variables $\rho,\kappa$ to the env
state; initialise to $0$ each episode. Replace the per-step $\reff=\rbase+\Delta
r(a)$, $\Keff=\Kbase+\Delta K(a)$ with the cumulative update of
Eq.~\eqref{eq:cumdyn}: accumulate $\Delta r,\Delta K$, then clip to
$[\rmin,\rmax]$, $[\Kmin,\Kmax]$.
\item Implement the caps of Section~\ref{sec:tier3}: $\rmax{=}0.55$ for
Ricker/Allee/regime; $\rmax(\theta){=}1.6/\theta$ for theta (the simulator knows
$\theta$); $\rmin{=}-0.15$; $\Kmin{=}500$, $\Kmax{=}1500$. Expose them in the
per-cell config. Add an optional decay $(d_r,d_K)$ defaulting to $0$.
\item Replace both action tables with the revised $5$- and $10$-action tables
(Section~\ref{sec:tier3}). Harvest is a negative $\Delta r$ with negative cost
(revenue). Remove any direct-$s$ authority ($h(a),\delta(a)$) if present in the
code --- Tier-3 is model-only.
\item Fix the theta prior to guarantee base-stability headroom: draw $\theta$ then
$\rbase\sim\mathcal{U}(0.12,\min(0.40,1.2/\theta))$ for theta cells.
\item Expose $(r_{\mathrm{eff},t},K_{\mathrm{eff},t})$ (equivalently $\rho_t,\kappa_t$) to every
method as observed inputs; only $s$ stays partially observed.
\item Keep the reward with a \emph{fixed} $\Kref{=}500$ (Eq.~\eqref{eq:reward3});
do \emph{not} switch the saturation reference to $\Keff$ (Section~\ref{sec:rewardfix}).
\end{enumerate}

\paragraph{(C1) Belief filter with parameter identification (Tier~A \#1).}
Extend the particle filter to joint $(s,m)$ inference: particles carry
$(s,\rbase,C,\theta,z)$ with the per-cell priors; use the learned dynamics as the
proposal and the exact log-normal emission as the weight. Prefer Rao--Blackwellization
(grid/analytic over the low-dimensional structural parameters, particles over $s$)
to avoid particle depletion in the higher-dimensional latent; raise the particle
budget only if a depletion diagnostic (ESS) demands it. Output a posterior over
$m$, not just $\hat s$. Add an \texttt{oracle\_state} and an \texttt{oracle\_params}
switch (feed true $s$ and/or true $m$) for the Tier~B ablation.

\paragraph{(C2) Belief-expected reward (Tier~A \#2).}
Add a reward mode that returns $r(b_t,a)=\mathbb{E}_{s\sim b_t}[\,\cdot\,]$ computed
from the current belief (expectation of the abundance term and of the collapse
indicator), used by every planner during rollout and by the learners' training
target. Keep the environment's realized $\Rtrue$ as the \emph{evaluation} channel
only; log both and report $\Rtrue$ for all cross-$\sobs$ comparisons.

\paragraph{(C3) Richer dynamics member (Tier~C \#8).}
Add a small-MLP ensemble member as an option alongside the ridge member, composed
with the exact emission (never trained on $o_{t+1}$ directly). Keep ensemble size,
MPC horizon, and pessimism $\lambda$ swappable.

\paragraph{(C4) Delphic as model-based (Tier~C \#6).}
Replace the random-feature-linear CQL head with: an ensemble of compatible world
models drawn from the filter's $m$-posterior; plan on the shared MPC backbone;
penalize the planned value by the cross-world variance $u_\Delta=\mathrm{Var}_w
Q^\pi_w$. Add a check that $u_\Delta$ is monotone in $\sobs$ and $\to 0$ at
$\sobs{=}0$; sweep $\beta_\Delta$.

\paragraph{(C5) OGSRL budget recalibration (Tier~C \#7).}
Separate two quantities that Tier-2 conflated: the training-time discounted
\emph{cumulative} safety/OOD budgets, and the deployment-time \emph{per-step}
feasibility threshold. Derive the cumulative budget from the target episode
collapse probability (use $\sum_t\gamma^t\approx 18.5$ to convert), and tune the
guardian quantile per $\sobs$ (it was too tight at $\sobs{=}0$). Keep the dual
ascent and least-violating fallback; log the fallback rate.

\paragraph{(C6) BA-MCTS safety-aware leaves (Tier~C \#8).}
Replace zero leaf values with a fitted belief-value (shared with RefPlan) and add a
near-threshold penalty to the leaf/expansion value.

\paragraph{(C7) Diagnostics and evaluation (Tier~B, Tier~D).}
\begin{enumerate}[leftmargin=1.8em,nosep]
\item Run the decision-relevance gate \emph{at every} $\sobs$ including $0.4$, and
record the reward/collapse gaps; treat $\sobs{=}0.4$ as a reported diagnostic.
\item Run the oracle-filter (and oracle-params) ablation at $\sobs{=}0.4$ across all
methods.
\item In the results extractor, report both MOOR numbers (shared learned filter and
own Ricker filter) and $\Rtrue$ alongside operational return; add collapse split by
start condition (above vs.\ below $C$).
\item Add independently seeded datasets per cell and report paired-bootstrap CIs on
the beats-both rate as the headline statistic.
\end{enumerate}

\paragraph{(C8) Recalibration (gate everything).}
After C0--C2, recalibrate the behavior-policy mixture so the healthy-start incident
collapse rate returns to $[0.15,0.24]$ under the new cumulative dynamics, and
re-pass the decision gate at $\sobs\!\le\!0.2$ before reading the high-noise cells.
The cumulative controls and $\rmin{<}0$ harvest change the achievable collapse
distribution, so the Tier-2 profiles will not transfer unchanged.

\paragraph{Suggested order.} C0 $\to$ C8 (env $+$ recalibrate, so the new setting is
sound) $\to$ C1, C2 (the Tier-A levers) $\to$ C7 oracle/gate diagnostics (decide
whether $\sobs{=}0.4$ is recoverable) $\to$ C3--C6 (per-method) $\to$ C7 statistics.

\begin{thebibliography}{9}
\bibitem{delphic2024} Pace, A., Y\`eche, H., Sch\"olkopf, B., R\"atsch, G.\ \& Tennenholtz, G.\ (2024).
Delphic Offline Reinforcement Learning under Nonidentifiable Hidden Confounding. \emph{ICLR}. arXiv:2306.01157.
\bibitem{ogsrl2025} Yan, R., Shen, X., Wachi, A., Gros, S., Zhao, A.\ \& Hu, X.\ (2025).
Offline Guarded Safe Reinforcement Learning for Medical Treatment Optimization Strategies. \emph{NeurIPS}. arXiv:2505.16242.
\bibitem{confpomdp2024} Hong, M., Qi, Z.\ \& Xu, Y.\ (2024).
Model-based Reinforcement Learning for Confounded POMDPs. \emph{ICML}, PMLR v235.
\bibitem{mopo2020} Yu, T.\ et al.\ (2020). MOPO: Model-based Offline Policy Optimization. \emph{NeurIPS}.
\bibitem{refplan2025} Reflect-then-Plan: Offline Model-Based Planning through a Doubly Bayesian Lens. \emph{ICML} (2025).
\bibitem{bamcts2026} Bayes-Adaptive Monte Carlo Tree Search for Offline Model-based RL. \emph{ICLR} (2026).
\bibitem{plus2018} Memarzadeh, M.\ \& Boettiger, C.\ (2018). Adaptive management of ecological systems under partial observability. \emph{Biological Conservation}, 224, 9--15.
\bibitem{moor2025} Ju, H.\ et al.\ (2025). Mechanistic model-based offline RL for fishery management. \emph{Expert Systems}.
\end{thebibliography}

\end{document}